% !TeX root = Master.tex
\section{Discussion}
\label{sec:discussion}

This paper formalizes the objectives and assumptions of empirical Monte Carlo benchmarking and examines whether fitted outcome simulators preserve estimator rankings. The results suggest that ranking fidelity depends on the simulator, the candidate estimators, and the underlying DGP. The variation across simulators is consistent with the sensitivity to simulation models reported by \authorcite{gentzelCaseEvaluatingCausal} and \authorcite{amaranathImprovingGenerativeMethods}. Variation across N\&W designs was greater than across ACIC error structures, pointing to the importance of conditional outcome structure and treatment confounding; RF under non-additive errors was the clearest error-structure exception. Larger reference samples sometimes substantially changed fidelity, but did not consistently improve it. Aggregating rankings may offer a robustness check under model uncertainty, although it never outperformed the best individual simulator and was not uniformly reliable.

Model alignment is a central concern because a simulator can favor estimators sharing its outcome-model specification, particularly when that model poorly approximates the target outcome structure. G-computation showed the strongest promotions, TMLE was less affected, and AIPW generally showed little promotion or slight demotion. Excluding aligned estimators substantially improved ranking fidelity in some settings, but poor fidelity sometimes persisted and exclusion did not always help. Alignment therefore explains part of the disagreement between benchmarks, rather than all simulator failures. Its model-dependence agrees with \authorcite{schulerSynthValidationSelectingBest2017}, who found alignment effects for linear regression and unconstrained gradient boosting, and \authorcite{karimPlasmodeBenchmarkCan2026}, who demonstrated alignment bias for a parametric outcome model. OLS2 showed the largest exclusion-related improvements in ACIC, while the effects for flexible learners were more varied.

Estimand fidelity and ranking fidelity were generally closely related. Nevertheless, similar mean absolute estimand differences sometimes accompanied substantially different rankings, even after excluding aligned estimators. These results provide some evidence that matching the reference effect is not enough to ensure ranking fidelity, illustrating the distinction between estimand transfer and ranking transfer made by \authorcite{karimPlasmodeBenchmarkCan2026}.

The regret comparisons answer the question of whether conducting an EMCS for estimator selection adds value over using the associated G-computation estimator. In line with \authorcite{advaniMostlyHarmlessSimulations2019}, misspecified models can mislead estimator comparisons. Our regret results also show that such simulators may still be useful for selection when aligned estimators are excluded, potentially reducing regret under model uncertainty. When an outcome model is well supported, direct G-computation may be preferable, while selection from its full induced ranking generally incurs a comparatively small penalty. The value of an EMCS thus depends on prior confidence in the model and which estimators remain eligible for selection.

The low ranking fidelity of WGAN in ACIC despite reasonable distributional diagnostics (not shown) suggests that checks similar to those proposed by \authorcite{atheyUsingWassersteinGenerative2024} are insufficient on their own to validate a benchmark for estimator selection. The distance associations and sensitivity checks suggest that outcome similarity may identify pronounced simulator failures, but can be less useful for distinguishing among the remaining alternatives. The greater robustness of the N\&W B--D associations shows that this limitation is setting-dependent.

Consequently, no simulator should be expected to provide universally reliable rankings across all empirical settings. A practical approach could be to fit several simulation models, use outcome distances as exploratory diagnostics alongside prior knowledge and other checks to assess confidence in each, and compare the resulting rankings after excluding aligned estimators. An aggregate ranking can provide a robustness check or a default when there is little basis for preferring one simulator, although our results do not establish it as superior to individual models. When an outcome model is already well supported, using it directly for G-comp or selecting an estimator from its full induced ranking, without excluding aligned estimators, may be preferable.

Finally, several limitations should be noted. Our simulation design holds covariates and treatment fixed within each benchmark and simulates only outcomes. The findings therefore only partially extend to EMCS that also model treatment assignment or the covariate distribution. Further work is needed to assess how approximation of these additional components, such as the propensity score, can distort estimator comparisons and interact with outcome-model approximation and model alignment. The evidence comes from two families of designs (ACIC 2017 and the N\&W designs with six covariates), with simulators and outcome learners fitted with fixed hyperparameter settings to facilitate comparisons and reduce variation from tuning. It remains unclear how well the findings generalize to other designs, or how much tuning would change the results or improve the performance of models such as RF. The WGAN results concern a single specification; alternative architectures and training procedures may yield different findings. Comparing alternative models and architectures could help establish whether the failures observed here persist and clarify how model alignment arises, including within model families.

Another limitation of our evaluation is our reliance mainly on averages across reference datasets, although we performed a sensitivity check using medians and obtained generally similar findings. The biggest differences concerned regret magnitudes, which sometimes changed substantially, although the mean-optimal selection rule remained optimal under the median in all cases. Our evaluation relies mainly on Spearman rank correlation and estimator-selection regret; although checks using Kendall's \(\tau\) gave largely consistent comparisons, other measures could reveal aspects of benchmark performance that these summaries miss. We do not show these robustness checks. Estimator comparison focused on MSE only. A particularly interesting direction for future work could be to extend evaluations to the accuracy of estimated standard errors and confidence-interval coverage. This would also assess whether simulators reproduce inferential performance, which preservation of MSE rankings alone does not guarantee.
